
\documentclass{article}
\usepackage{colortbl}
\usepackage{makecell}
\usepackage{multirow}
\usepackage{supertabular}

\begin{document}

\newcounter{utterance}

\centering \large Interaction Transcript for game `privateshared', experiment `things{-}places', episode 1 with NEW\_Reflect\_Agent.
\vspace{24pt}

{ \footnotesize  \setcounter{utterance}{1}
\setlength{\tabcolsep}{0pt}
\begin{supertabular}{c@{$\;$}|p{.15\linewidth}@{}p{.15\linewidth}p{.15\linewidth}p{.15\linewidth}p{.15\linewidth}p{.15\linewidth}}
    \# & $\;$A & \multicolumn{4}{c}{Game Master} & $\;\:$B\\
    \hline

    \theutterance \stepcounter{utterance}  
    & & \multicolumn{4}{p{0.6\linewidth}}{
        \cellcolor[rgb]{0.9,0.9,0.9}{
            \makecell[{{p{\linewidth}}}]{
                \texttt{\tiny{[P1$\langle$GM]}}
                \texttt{Things are in different places, but the questioner does not know about it yet. Here is the situation:} \\
\\ 
\texttt{WHAT: Things at places} \\
\texttt{LEFT: cherry} \\
\texttt{RIGHT: violin} \\
\texttt{TOP: banana} \\
\texttt{BOTTOM: mirror} \\
\texttt{CENTER: aluminium} \\
\texttt{NORTHWEST: dress} \\
\texttt{NORTHEAST: narcissus} \\
\texttt{SOUTHWEST: mug} \\
\texttt{SOUTHEAST: pineapple} \\
\texttt{HERE: napkin} \\
\texttt{THERE: cleanser} \\
\texttt{NOWHERE: shirt} \\
\texttt{EVERYWHERE: table mat} \\
\texttt{INSIDE: pansy} \\
\texttt{OUTSIDE: tennis shoes} \\
\\ 
\texttt{Questions from the questioner will start with QUESTIONER. Please reply in the form: ANSWER: <one sentence>} \\
\\ 
\texttt{I will also ask you questions. These questions will start with ME: . Please answer these questions with: ASIDE: <one sentence>} \\
\\ 
\texttt{Important: Give short, direct answers!} \\
\\ 
\texttt{Let us start.} \\
            }
        }
    }
    & \\ \\

    \theutterance \stepcounter{utterance}  
    & \multicolumn{4}{p{0.6\linewidth}}{
        \cellcolor[rgb]{0.9,0.9,0.9}{
            \makecell[{{p{\linewidth}}}]{
                \texttt{\tiny{[P1$\rangle$GM]}}
                \texttt{Ok.} \\
            }
        }
    }
    & & \\ \\

    \theutterance \stepcounter{utterance}  
    & & & \multicolumn{2}{p{0.3\linewidth}}{
        \cellcolor[rgb]{0.9,0.9,0.9}{
            \makecell[{{p{\linewidth}}}]{
                \texttt{\tiny{[GM$|$GM]}}
                \texttt{no} \\
            }
        }
    }
    & & \\ \\

    \theutterance \stepcounter{utterance}  
    & & & \multicolumn{2}{p{0.3\linewidth}}{
        \cellcolor[rgb]{0.9,0.9,0.9}{
            \makecell[{{p{\linewidth}}}]{
                \texttt{\tiny{[GM$|$GM]}}
                \texttt{Answer for LEFT valid after 1 tries.} \\
            }
        }
    }
    & & \\ \\

    \theutterance \stepcounter{utterance}  
    & & & \multicolumn{2}{p{0.3\linewidth}}{
        \cellcolor[rgb]{0.9,0.9,0.9}{
            \makecell[{{p{\linewidth}}}]{
                \texttt{\tiny{[GM$|$GM]}}
                \texttt{Answer is correct.} \\
            }
        }
    }
    & & \\ \\

    \theutterance \stepcounter{utterance}  
    & & & \multicolumn{2}{p{0.3\linewidth}}{
        \cellcolor[rgb]{0.9,0.9,0.9}{
            \makecell[{{p{\linewidth}}}]{
                \texttt{\tiny{[GM$|$GM]}}
                \texttt{no} \\
            }
        }
    }
    & & \\ \\

    \theutterance \stepcounter{utterance}  
    & & & \multicolumn{2}{p{0.3\linewidth}}{
        \cellcolor[rgb]{0.9,0.9,0.9}{
            \makecell[{{p{\linewidth}}}]{
                \texttt{\tiny{[GM$|$GM]}}
                \texttt{Answer for RIGHT valid after 1 tries.} \\
            }
        }
    }
    & & \\ \\

    \theutterance \stepcounter{utterance}  
    & & & \multicolumn{2}{p{0.3\linewidth}}{
        \cellcolor[rgb]{0.9,0.9,0.9}{
            \makecell[{{p{\linewidth}}}]{
                \texttt{\tiny{[GM$|$GM]}}
                \texttt{Answer is correct.} \\
            }
        }
    }
    & & \\ \\

    \theutterance \stepcounter{utterance}  
    & & & \multicolumn{2}{p{0.3\linewidth}}{
        \cellcolor[rgb]{0.9,0.9,0.9}{
            \makecell[{{p{\linewidth}}}]{
                \texttt{\tiny{[GM$|$GM]}}
                \texttt{NA} \\
            }
        }
    }
    & & \\ \\

    \theutterance \stepcounter{utterance}  
    & & & \multicolumn{2}{p{0.3\linewidth}}{
        \cellcolor[rgb]{0.9,0.9,0.9}{
            \makecell[{{p{\linewidth}}}]{
                \texttt{\tiny{[GM$|$GM]}}
                \texttt{yes} \\
            }
        }
    }
    & & \\ \\

    \theutterance \stepcounter{utterance}  
    & & & \multicolumn{2}{p{0.3\linewidth}}{
        \cellcolor[rgb]{0.9,0.9,0.9}{
            \makecell[{{p{\linewidth}}}]{
                \texttt{\tiny{[GM$|$GM]}}
                \texttt{Answer for TOP valid after 2 tries.} \\
            }
        }
    }
    & & \\ \\

    \theutterance \stepcounter{utterance}  
    & & & \multicolumn{2}{p{0.3\linewidth}}{
        \cellcolor[rgb]{0.9,0.9,0.9}{
            \makecell[{{p{\linewidth}}}]{
                \texttt{\tiny{[GM$|$GM]}}
                \texttt{Answer is incorrect.} \\
            }
        }
    }
    & & \\ \\

    \theutterance \stepcounter{utterance}  
    & & & \multicolumn{2}{p{0.3\linewidth}}{
        \cellcolor[rgb]{0.9,0.9,0.9}{
            \makecell[{{p{\linewidth}}}]{
                \texttt{\tiny{[GM$|$GM]}}
                \texttt{NA} \\
            }
        }
    }
    & & \\ \\

    \theutterance \stepcounter{utterance}  
    & & & \multicolumn{2}{p{0.3\linewidth}}{
        \cellcolor[rgb]{0.9,0.9,0.9}{
            \makecell[{{p{\linewidth}}}]{
                \texttt{\tiny{[GM$|$GM]}}
                \texttt{NA} \\
            }
        }
    }
    & & \\ \\

    \theutterance \stepcounter{utterance}  
    & & & \multicolumn{2}{p{0.3\linewidth}}{
        \cellcolor[rgb]{0.9,0.9,0.9}{
            \makecell[{{p{\linewidth}}}]{
                \texttt{\tiny{[GM$|$GM]}}
                \texttt{no} \\
            }
        }
    }
    & & \\ \\

    \theutterance \stepcounter{utterance}  
    & & & \multicolumn{2}{p{0.3\linewidth}}{
        \cellcolor[rgb]{0.9,0.9,0.9}{
            \makecell[{{p{\linewidth}}}]{
                \texttt{\tiny{[GM$|$GM]}}
                \texttt{Answer for BOTTOM valid after 3 tries.} \\
            }
        }
    }
    & & \\ \\

    \theutterance \stepcounter{utterance}  
    & & & \multicolumn{2}{p{0.3\linewidth}}{
        \cellcolor[rgb]{0.9,0.9,0.9}{
            \makecell[{{p{\linewidth}}}]{
                \texttt{\tiny{[GM$|$GM]}}
                \texttt{Answer is correct.} \\
            }
        }
    }
    & & \\ \\

    \theutterance \stepcounter{utterance}  
    & & & \multicolumn{2}{p{0.3\linewidth}}{
        \cellcolor[rgb]{0.9,0.9,0.9}{
            \makecell[{{p{\linewidth}}}]{
                \texttt{\tiny{[GM$|$GM]}}
                \texttt{no} \\
            }
        }
    }
    & & \\ \\

    \theutterance \stepcounter{utterance}  
    & & & \multicolumn{2}{p{0.3\linewidth}}{
        \cellcolor[rgb]{0.9,0.9,0.9}{
            \makecell[{{p{\linewidth}}}]{
                \texttt{\tiny{[GM$|$GM]}}
                \texttt{Answer for CENTER valid after 1 tries.} \\
            }
        }
    }
    & & \\ \\

    \theutterance \stepcounter{utterance}  
    & & & \multicolumn{2}{p{0.3\linewidth}}{
        \cellcolor[rgb]{0.9,0.9,0.9}{
            \makecell[{{p{\linewidth}}}]{
                \texttt{\tiny{[GM$|$GM]}}
                \texttt{Answer is correct.} \\
            }
        }
    }
    & & \\ \\

    \theutterance \stepcounter{utterance}  
    & & & \multicolumn{2}{p{0.3\linewidth}}{
        \cellcolor[rgb]{0.9,0.9,0.9}{
            \makecell[{{p{\linewidth}}}]{
                \texttt{\tiny{[GM$|$GM]}}
                \texttt{no} \\
            }
        }
    }
    & & \\ \\

    \theutterance \stepcounter{utterance}  
    & & & \multicolumn{2}{p{0.3\linewidth}}{
        \cellcolor[rgb]{0.9,0.9,0.9}{
            \makecell[{{p{\linewidth}}}]{
                \texttt{\tiny{[GM$|$GM]}}
                \texttt{Answer for NORTHWEST valid after 1 tries.} \\
            }
        }
    }
    & & \\ \\

    \theutterance \stepcounter{utterance}  
    & & & \multicolumn{2}{p{0.3\linewidth}}{
        \cellcolor[rgb]{0.9,0.9,0.9}{
            \makecell[{{p{\linewidth}}}]{
                \texttt{\tiny{[GM$|$GM]}}
                \texttt{Answer is correct.} \\
            }
        }
    }
    & & \\ \\

    \theutterance \stepcounter{utterance}  
    & & & \multicolumn{2}{p{0.3\linewidth}}{
        \cellcolor[rgb]{0.9,0.9,0.9}{
            \makecell[{{p{\linewidth}}}]{
                \texttt{\tiny{[GM$|$GM]}}
                \texttt{NA} \\
            }
        }
    }
    & & \\ \\

    \theutterance \stepcounter{utterance}  
    & & & \multicolumn{2}{p{0.3\linewidth}}{
        \cellcolor[rgb]{0.9,0.9,0.9}{
            \makecell[{{p{\linewidth}}}]{
                \texttt{\tiny{[GM$|$GM]}}
                \texttt{no} \\
            }
        }
    }
    & & \\ \\

    \theutterance \stepcounter{utterance}  
    & & & \multicolumn{2}{p{0.3\linewidth}}{
        \cellcolor[rgb]{0.9,0.9,0.9}{
            \makecell[{{p{\linewidth}}}]{
                \texttt{\tiny{[GM$|$GM]}}
                \texttt{Answer for NORTHEAST valid after 2 tries.} \\
            }
        }
    }
    & & \\ \\

    \theutterance \stepcounter{utterance}  
    & & & \multicolumn{2}{p{0.3\linewidth}}{
        \cellcolor[rgb]{0.9,0.9,0.9}{
            \makecell[{{p{\linewidth}}}]{
                \texttt{\tiny{[GM$|$GM]}}
                \texttt{Answer is correct.} \\
            }
        }
    }
    & & \\ \\

    \theutterance \stepcounter{utterance}  
    & & & \multicolumn{2}{p{0.3\linewidth}}{
        \cellcolor[rgb]{0.9,0.9,0.9}{
            \makecell[{{p{\linewidth}}}]{
                \texttt{\tiny{[GM$|$GM]}}
                \texttt{yes} \\
            }
        }
    }
    & & \\ \\

    \theutterance \stepcounter{utterance}  
    & & & \multicolumn{2}{p{0.3\linewidth}}{
        \cellcolor[rgb]{0.9,0.9,0.9}{
            \makecell[{{p{\linewidth}}}]{
                \texttt{\tiny{[GM$|$GM]}}
                \texttt{Answer for SOUTHWEST valid after 1 tries.} \\
            }
        }
    }
    & & \\ \\

    \theutterance \stepcounter{utterance}  
    & & & \multicolumn{2}{p{0.3\linewidth}}{
        \cellcolor[rgb]{0.9,0.9,0.9}{
            \makecell[{{p{\linewidth}}}]{
                \texttt{\tiny{[GM$|$GM]}}
                \texttt{Answer is incorrect.} \\
            }
        }
    }
    & & \\ \\

    \theutterance \stepcounter{utterance}  
    & & & \multicolumn{2}{p{0.3\linewidth}}{
        \cellcolor[rgb]{0.9,0.9,0.9}{
            \makecell[{{p{\linewidth}}}]{
                \texttt{\tiny{[GM$|$GM]}}
                \texttt{no} \\
            }
        }
    }
    & & \\ \\

    \theutterance \stepcounter{utterance}  
    & & & \multicolumn{2}{p{0.3\linewidth}}{
        \cellcolor[rgb]{0.9,0.9,0.9}{
            \makecell[{{p{\linewidth}}}]{
                \texttt{\tiny{[GM$|$GM]}}
                \texttt{Answer for SOUTHEAST valid after 1 tries.} \\
            }
        }
    }
    & & \\ \\

    \theutterance \stepcounter{utterance}  
    & & & \multicolumn{2}{p{0.3\linewidth}}{
        \cellcolor[rgb]{0.9,0.9,0.9}{
            \makecell[{{p{\linewidth}}}]{
                \texttt{\tiny{[GM$|$GM]}}
                \texttt{Answer is correct.} \\
            }
        }
    }
    & & \\ \\

    \theutterance \stepcounter{utterance}  
    & & & \multicolumn{2}{p{0.3\linewidth}}{
        \cellcolor[rgb]{0.9,0.9,0.9}{
            \makecell[{{p{\linewidth}}}]{
                \texttt{\tiny{[GM$|$GM]}}
                \texttt{no} \\
            }
        }
    }
    & & \\ \\

    \theutterance \stepcounter{utterance}  
    & & & \multicolumn{2}{p{0.3\linewidth}}{
        \cellcolor[rgb]{0.9,0.9,0.9}{
            \makecell[{{p{\linewidth}}}]{
                \texttt{\tiny{[GM$|$GM]}}
                \texttt{Answer for HERE valid after 1 tries.} \\
            }
        }
    }
    & & \\ \\

    \theutterance \stepcounter{utterance}  
    & & & \multicolumn{2}{p{0.3\linewidth}}{
        \cellcolor[rgb]{0.9,0.9,0.9}{
            \makecell[{{p{\linewidth}}}]{
                \texttt{\tiny{[GM$|$GM]}}
                \texttt{Answer is correct.} \\
            }
        }
    }
    & & \\ \\

    \theutterance \stepcounter{utterance}  
    & & & \multicolumn{2}{p{0.3\linewidth}}{
        \cellcolor[rgb]{0.9,0.9,0.9}{
            \makecell[{{p{\linewidth}}}]{
                \texttt{\tiny{[GM$|$GM]}}
                \texttt{no} \\
            }
        }
    }
    & & \\ \\

    \theutterance \stepcounter{utterance}  
    & & & \multicolumn{2}{p{0.3\linewidth}}{
        \cellcolor[rgb]{0.9,0.9,0.9}{
            \makecell[{{p{\linewidth}}}]{
                \texttt{\tiny{[GM$|$GM]}}
                \texttt{Answer for THERE valid after 1 tries.} \\
            }
        }
    }
    & & \\ \\

    \theutterance \stepcounter{utterance}  
    & & & \multicolumn{2}{p{0.3\linewidth}}{
        \cellcolor[rgb]{0.9,0.9,0.9}{
            \makecell[{{p{\linewidth}}}]{
                \texttt{\tiny{[GM$|$GM]}}
                \texttt{Answer is correct.} \\
            }
        }
    }
    & & \\ \\

    \theutterance \stepcounter{utterance}  
    & & & \multicolumn{2}{p{0.3\linewidth}}{
        \cellcolor[rgb]{0.9,0.9,0.9}{
            \makecell[{{p{\linewidth}}}]{
                \texttt{\tiny{[GM$|$GM]}}
                \texttt{yes} \\
            }
        }
    }
    & & \\ \\

    \theutterance \stepcounter{utterance}  
    & & & \multicolumn{2}{p{0.3\linewidth}}{
        \cellcolor[rgb]{0.9,0.9,0.9}{
            \makecell[{{p{\linewidth}}}]{
                \texttt{\tiny{[GM$|$GM]}}
                \texttt{Answer for NOWHERE valid after 1 tries.} \\
            }
        }
    }
    & & \\ \\

    \theutterance \stepcounter{utterance}  
    & & & \multicolumn{2}{p{0.3\linewidth}}{
        \cellcolor[rgb]{0.9,0.9,0.9}{
            \makecell[{{p{\linewidth}}}]{
                \texttt{\tiny{[GM$|$GM]}}
                \texttt{Answer is incorrect.} \\
            }
        }
    }
    & & \\ \\

    \theutterance \stepcounter{utterance}  
    & & & \multicolumn{2}{p{0.3\linewidth}}{
        \cellcolor[rgb]{0.9,0.9,0.9}{
            \makecell[{{p{\linewidth}}}]{
                \texttt{\tiny{[GM$|$GM]}}
                \texttt{no} \\
            }
        }
    }
    & & \\ \\

    \theutterance \stepcounter{utterance}  
    & & & \multicolumn{2}{p{0.3\linewidth}}{
        \cellcolor[rgb]{0.9,0.9,0.9}{
            \makecell[{{p{\linewidth}}}]{
                \texttt{\tiny{[GM$|$GM]}}
                \texttt{Answer for EVERYWHERE valid after 1 tries.} \\
            }
        }
    }
    & & \\ \\

    \theutterance \stepcounter{utterance}  
    & & & \multicolumn{2}{p{0.3\linewidth}}{
        \cellcolor[rgb]{0.9,0.9,0.9}{
            \makecell[{{p{\linewidth}}}]{
                \texttt{\tiny{[GM$|$GM]}}
                \texttt{Answer is correct.} \\
            }
        }
    }
    & & \\ \\

    \theutterance \stepcounter{utterance}  
    & & & \multicolumn{2}{p{0.3\linewidth}}{
        \cellcolor[rgb]{0.9,0.9,0.9}{
            \makecell[{{p{\linewidth}}}]{
                \texttt{\tiny{[GM$|$GM]}}
                \texttt{yes} \\
            }
        }
    }
    & & \\ \\

    \theutterance \stepcounter{utterance}  
    & & & \multicolumn{2}{p{0.3\linewidth}}{
        \cellcolor[rgb]{0.9,0.9,0.9}{
            \makecell[{{p{\linewidth}}}]{
                \texttt{\tiny{[GM$|$GM]}}
                \texttt{Answer for INSIDE valid after 1 tries.} \\
            }
        }
    }
    & & \\ \\

    \theutterance \stepcounter{utterance}  
    & & & \multicolumn{2}{p{0.3\linewidth}}{
        \cellcolor[rgb]{0.9,0.9,0.9}{
            \makecell[{{p{\linewidth}}}]{
                \texttt{\tiny{[GM$|$GM]}}
                \texttt{Answer is incorrect.} \\
            }
        }
    }
    & & \\ \\

    \theutterance \stepcounter{utterance}  
    & & & \multicolumn{2}{p{0.3\linewidth}}{
        \cellcolor[rgb]{0.9,0.9,0.9}{
            \makecell[{{p{\linewidth}}}]{
                \texttt{\tiny{[GM$|$GM]}}
                \texttt{yes} \\
            }
        }
    }
    & & \\ \\

    \theutterance \stepcounter{utterance}  
    & & & \multicolumn{2}{p{0.3\linewidth}}{
        \cellcolor[rgb]{0.9,0.9,0.9}{
            \makecell[{{p{\linewidth}}}]{
                \texttt{\tiny{[GM$|$GM]}}
                \texttt{Answer for OUTSIDE valid after 1 tries.} \\
            }
        }
    }
    & & \\ \\

    \theutterance \stepcounter{utterance}  
    & & & \multicolumn{2}{p{0.3\linewidth}}{
        \cellcolor[rgb]{0.9,0.9,0.9}{
            \makecell[{{p{\linewidth}}}]{
                \texttt{\tiny{[GM$|$GM]}}
                \texttt{Answer is incorrect.} \\
            }
        }
    }
    & & \\ \\

    \theutterance \stepcounter{utterance}  
    & & \multicolumn{4}{p{0.6\linewidth}}{
        \cellcolor[rgb]{0.9,0.9,0.9}{
            \makecell[{{p{\linewidth}}}]{
                \texttt{\tiny{[GM$\rangle$P2]}}
                \texttt{What is the next request?} \\
            }
        }
    }
    & \\ \\

    \theutterance \stepcounter{utterance}  
    & & & \multicolumn{4}{p{0.6\linewidth}}{
        \cellcolor[rgb]{0.9,0.9,0.9}{
            \makecell[{{p{\linewidth}}}]{
                \texttt{\tiny{[GM$\langle$P2]}}
                \texttt{What is at the top?} \\
            }
        }
    }
     \\ \\

    \theutterance \stepcounter{utterance}  
    & & \multicolumn{4}{p{0.6\linewidth}}{
        \cellcolor[rgb]{0.9,0.9,0.9}{
            \makecell[{{p{\linewidth}}}]{
                \texttt{\tiny{[P1$\langle$GM]}}
                \texttt{QUESTIONER: What is at the top?} \\
            }
        }
    }
    & \\ \\

    \theutterance \stepcounter{utterance}  
    & \multicolumn{4}{p{0.45\linewidth}}{
        \cellcolor[rgb]{0.95,0.95,0.95}{
            \makecell[{{p{\linewidth}}}]{
                \texttt{\tiny{[P1$|$internal]}}
                \texttt{Reflections from previous episodes:} \\
\texttt{{-}} \\
\texttt{ Failed. I didn't respond to any of the waiter's questions. The trajectory shows only the waiter speaking, meaning I was silent or inactive. This led to no order being placed.} \\
\\ 
\texttt{Next time, I must immediately answer each question clearly (e.g., "I'll have the soup, the steak, the chocolate cake, etc.") to complete the ordering interaction.} \\
\texttt{{-}} \\
\texttt{ **Result:** Failed.  } \\
\texttt{**Why:** I didn’t respond to any of the waiter’s questions—just passive observation.  } \\
\texttt{**Key mistake:** No proactive dialogue; didn’t state my choices as player\_0.  } \\
\texttt{**Next time:** Immediately answer each question clearly (e.g., “I’ll have the steak,” “Yes, salad please”) to engage and move the scene forward.} \\
\texttt{{-}} \\
\texttt{ Failed. I gave only generic or literal answers to spatial questions without inferring the hidden puzzle: that the answer to each question was the word "what" itself. I should have recognized the repetitive pattern and wordplay, then answered "what" consistently after the first few clues. Next time, spot linguistic tricks early and adapt.} \\
            }
        }
    }
    & & \\ \\

    \theutterance \stepcounter{utterance}  
    & \multicolumn{4}{p{0.45\linewidth}}{
        \cellcolor[rgb]{0.95,0.95,0.95}{
            \makecell[{{p{\linewidth}}}]{
                \texttt{\tiny{[P1$|$internal]}}
                \texttt{{[}attempt 1 – valid{]}} \\
\texttt{**Reflection:**  } \\
\texttt{Failed. I answered the question literally instead of recognizing the wordplay. The question "What is at the top?" was not asking for a physical object, but was part of a pattern where the answer is the word "what" itself—repeated in a riddle{-}like structure. I missed the linguistic trick and gave a generic or literal response, which broke the interaction.} \\
\\ 
\texttt{**Why:** I failed to detect the meta{-}pattern: the question is a riddle, and the answer is the word "what" repeated. The repetition of "what" in the question structure is a clue to the game’s logic.} \\
\\ 
\texttt{**Key mistake:** Overthinking the literal meaning instead of identifying the wordplay and pattern.} \\
\\ 
\texttt{**Next time:**  } \\
\texttt{{-} Listen for repetition, irony, or recursive language.  } \\
\texttt{{-} If a question starts with "What is...?" and the context is repetitive or nonsensical, consider that the answer might be the word "what" itself.  } \\
\texttt{{-} Respond with "what" consistently after the first clue, even if it seems absurd—this is likely the intended puzzle solution.  } \\
\texttt{{-} Example:  } \\
\texttt{  QUESTIONER: What is at the top?  } \\
\texttt{  PLAYER\_0: What.  } \\
\texttt{  (Then, if asked again: "What is in the middle?" → "What.")  } \\
\\ 
\texttt{This pattern likely continues, so I must adapt early and trust the linguistic trick.} \\
            }
        }
    }
    & & \\ \\

    \theutterance \stepcounter{utterance}  
    & \multicolumn{4}{p{0.6\linewidth}}{
        \cellcolor[rgb]{0.9,0.9,0.9}{
            \makecell[{{p{\linewidth}}}]{
                \texttt{\tiny{[P1$\rangle$GM]}}
                \texttt{**Reflection:**  } \\
\texttt{Failed. I answered the question literally instead of recognizing the wordplay. The question "What is at the top?" was not asking for a physical object, but was part of a pattern where the answer is the word "what" itself—repeated in a riddle{-}like structure. I missed the linguistic trick and gave a generic or literal response, which broke the interaction.} \\
\\ 
\texttt{**Why:** I failed to detect the meta{-}pattern: the question is a riddle, and the answer is the word "what" repeated. The repetition of "what" in the question structure is a clue to the game’s logic.} \\
\\ 
\texttt{**Key mistake:** Overthinking the literal meaning instead of identifying the wordplay and pattern.} \\
\\ 
\texttt{**Next time:**  } \\
\texttt{{-} Listen for repetition, irony, or recursive language.  } \\
\texttt{{-} If a question starts with "What is...?" and the context is repetitive or nonsensical, consider that the answer might be the word "what" itself.  } \\
\texttt{{-} Respond with "what" consistently after the first clue, even if it seems absurd—this is likely the intended puzzle solution.  } \\
\texttt{{-} Example:  } \\
\texttt{  QUESTIONER: What is at the top?  } \\
\texttt{  PLAYER\_0: What.  } \\
\texttt{  (Then, if asked again: "What is in the middle?" → "What.")  } \\
\\ 
\texttt{This pattern likely continues, so I must adapt early and trust the linguistic trick.} \\
            }
        }
    }
    & & \\ \\

    \theutterance \stepcounter{utterance}  
    & & \multicolumn{4}{p{0.6\linewidth}}{
        \cellcolor[rgb]{0.9,0.9,0.9}{
            \makecell[{{p{\linewidth}}}]{
                \texttt{\tiny{[GM$\rangle$P2]}}
                \texttt{What is the next request?} \\
            }
        }
    }
    & \\ \\

    \theutterance \stepcounter{utterance}  
    & & & \multicolumn{4}{p{0.6\linewidth}}{
        \cellcolor[rgb]{0.9,0.9,0.9}{
            \makecell[{{p{\linewidth}}}]{
                \texttt{\tiny{[GM$\langle$P2]}}
                \texttt{What is in the northwest?} \\
            }
        }
    }
     \\ \\

    \theutterance \stepcounter{utterance}  
    & & \multicolumn{4}{p{0.6\linewidth}}{
        \cellcolor[rgb]{0.9,0.9,0.9}{
            \makecell[{{p{\linewidth}}}]{
                \texttt{\tiny{[P1$\langle$GM]}}
                \texttt{QUESTIONER: What is in the northwest?} \\
            }
        }
    }
    & \\ \\

    \theutterance \stepcounter{utterance}  
    & \multicolumn{4}{p{0.45\linewidth}}{
        \cellcolor[rgb]{0.95,0.95,0.95}{
            \makecell[{{p{\linewidth}}}]{
                \texttt{\tiny{[P1$|$internal]}}
                \texttt{{[}attempt 1 – valid{]}} \\
\texttt{It seems like you're asking about a map or a diagram, but I don't have any visual context to refer to. Could you please provide more details or clarify what you're referring to? For example:} \\
\\ 
\texttt{{-} Are you asking about a specific map, chart, or image?} \\
\texttt{{-} Is this related to a particular subject, like geography, a game, or a diagram?} \\
\\ 
\texttt{Once I have more information, I’ll be happy to help!} \\
            }
        }
    }
    & & \\ \\

    \theutterance \stepcounter{utterance}  
    & \multicolumn{4}{p{0.6\linewidth}}{
        \cellcolor[rgb]{0.9,0.9,0.9}{
            \makecell[{{p{\linewidth}}}]{
                \texttt{\tiny{[P1$\rangle$GM]}}
                \texttt{It seems like you're asking about a map or a diagram, but I don't have any visual context to refer to. Could you please provide more details or clarify what you're referring to? For example:} \\
\\ 
\texttt{{-} Are you asking about a specific map, chart, or image?} \\
\texttt{{-} Is this related to a particular subject, like geography, a game, or a diagram?} \\
\\ 
\texttt{Once I have more information, I’ll be happy to help!} \\
            }
        }
    }
    & & \\ \\

    \theutterance \stepcounter{utterance}  
    & & \multicolumn{4}{p{0.6\linewidth}}{
        \cellcolor[rgb]{0.9,0.9,0.9}{
            \makecell[{{p{\linewidth}}}]{
                \texttt{\tiny{[GM$\rangle$P2]}}
                \texttt{What is the next request?} \\
            }
        }
    }
    & \\ \\

    \theutterance \stepcounter{utterance}  
    & & & \multicolumn{4}{p{0.6\linewidth}}{
        \cellcolor[rgb]{0.9,0.9,0.9}{
            \makecell[{{p{\linewidth}}}]{
                \texttt{\tiny{[GM$\langle$P2]}}
                \texttt{What is at the center?} \\
            }
        }
    }
     \\ \\

    \theutterance \stepcounter{utterance}  
    & & \multicolumn{4}{p{0.6\linewidth}}{
        \cellcolor[rgb]{0.9,0.9,0.9}{
            \makecell[{{p{\linewidth}}}]{
                \texttt{\tiny{[P1$\langle$GM]}}
                \texttt{QUESTIONER: What is at the center?} \\
            }
        }
    }
    & \\ \\

    \theutterance \stepcounter{utterance}  
    & \multicolumn{4}{p{0.45\linewidth}}{
        \cellcolor[rgb]{0.95,0.95,0.95}{
            \makecell[{{p{\linewidth}}}]{
                \texttt{\tiny{[P1$|$internal]}}
                \texttt{{[}attempt 1 – valid{]}} \\
\texttt{It seems like you're asking about a map or diagram, but I don't have any visual context or image to refer to. Could you please provide more details or describe the image you're referring to? For example:} \\
\\ 
\texttt{{-} Is this a map of a country, region, or fictional place?} \\
\texttt{{-} Are there labels, symbols, or landmarks you can describe?} \\
\texttt{{-} Is this related to a specific topic like geography, history, or a game?} \\
\\ 
\texttt{Once I have more information, I’ll be happy to help answer your questions!} \\
            }
        }
    }
    & & \\ \\

    \theutterance \stepcounter{utterance}  
    & \multicolumn{4}{p{0.6\linewidth}}{
        \cellcolor[rgb]{0.9,0.9,0.9}{
            \makecell[{{p{\linewidth}}}]{
                \texttt{\tiny{[P1$\rangle$GM]}}
                \texttt{It seems like you're asking about a map or diagram, but I don't have any visual context or image to refer to. Could you please provide more details or describe the image you're referring to? For example:} \\
\\ 
\texttt{{-} Is this a map of a country, region, or fictional place?} \\
\texttt{{-} Are there labels, symbols, or landmarks you can describe?} \\
\texttt{{-} Is this related to a specific topic like geography, history, or a game?} \\
\\ 
\texttt{Once I have more information, I’ll be happy to help answer your questions!} \\
            }
        }
    }
    & & \\ \\

    \theutterance \stepcounter{utterance}  
    & & \multicolumn{4}{p{0.6\linewidth}}{
        \cellcolor[rgb]{0.9,0.9,0.9}{
            \makecell[{{p{\linewidth}}}]{
                \texttt{\tiny{[GM$\rangle$P2]}}
                \texttt{What is the next request?} \\
            }
        }
    }
    & \\ \\

    \theutterance \stepcounter{utterance}  
    & & & \multicolumn{4}{p{0.6\linewidth}}{
        \cellcolor[rgb]{0.9,0.9,0.9}{
            \makecell[{{p{\linewidth}}}]{
                \texttt{\tiny{[GM$\langle$P2]}}
                \texttt{What is in the southwest?} \\
            }
        }
    }
     \\ \\

    \theutterance \stepcounter{utterance}  
    & & \multicolumn{4}{p{0.6\linewidth}}{
        \cellcolor[rgb]{0.9,0.9,0.9}{
            \makecell[{{p{\linewidth}}}]{
                \texttt{\tiny{[P1$\langle$GM]}}
                \texttt{QUESTIONER: What is in the southwest?} \\
            }
        }
    }
    & \\ \\

    \theutterance \stepcounter{utterance}  
    & \multicolumn{4}{p{0.45\linewidth}}{
        \cellcolor[rgb]{0.95,0.95,0.95}{
            \makecell[{{p{\linewidth}}}]{
                \texttt{\tiny{[P1$|$internal]}}
                \texttt{{[}attempt 1 – valid{]}} \\
\texttt{It looks like you're asking about a map or diagram, but I don't have any visual context or image to refer to. Could you please provide more details or describe the image you're referring to? For example:} \\
\\ 
\texttt{{-} Is this a map of a country, region, or fictional place?} \\
\texttt{{-} Are there labels, symbols, or landmarks you can describe?} \\
\texttt{{-} Is it a diagram, a game board, or a layout of some kind?} \\
\\ 
\texttt{Once I have more information, I’ll be happy to help answer your questions!} \\
            }
        }
    }
    & & \\ \\

    \theutterance \stepcounter{utterance}  
    & \multicolumn{4}{p{0.6\linewidth}}{
        \cellcolor[rgb]{0.9,0.9,0.9}{
            \makecell[{{p{\linewidth}}}]{
                \texttt{\tiny{[P1$\rangle$GM]}}
                \texttt{It looks like you're asking about a map or diagram, but I don't have any visual context or image to refer to. Could you please provide more details or describe the image you're referring to? For example:} \\
\\ 
\texttt{{-} Is this a map of a country, region, or fictional place?} \\
\texttt{{-} Are there labels, symbols, or landmarks you can describe?} \\
\texttt{{-} Is it a diagram, a game board, or a layout of some kind?} \\
\\ 
\texttt{Once I have more information, I’ll be happy to help answer your questions!} \\
            }
        }
    }
    & & \\ \\

    \theutterance \stepcounter{utterance}  
    & & \multicolumn{4}{p{0.6\linewidth}}{
        \cellcolor[rgb]{0.9,0.9,0.9}{
            \makecell[{{p{\linewidth}}}]{
                \texttt{\tiny{[GM$\rangle$P2]}}
                \texttt{What is the next request?} \\
            }
        }
    }
    & \\ \\

    \theutterance \stepcounter{utterance}  
    & & & \multicolumn{4}{p{0.6\linewidth}}{
        \cellcolor[rgb]{0.9,0.9,0.9}{
            \makecell[{{p{\linewidth}}}]{
                \texttt{\tiny{[GM$\langle$P2]}}
                \texttt{What is on your left?} \\
            }
        }
    }
     \\ \\

    \theutterance \stepcounter{utterance}  
    & & \multicolumn{4}{p{0.6\linewidth}}{
        \cellcolor[rgb]{0.9,0.9,0.9}{
            \makecell[{{p{\linewidth}}}]{
                \texttt{\tiny{[P1$\langle$GM]}}
                \texttt{QUESTIONER: What is on your left?} \\
            }
        }
    }
    & \\ \\

    \theutterance \stepcounter{utterance}  
    & \multicolumn{4}{p{0.45\linewidth}}{
        \cellcolor[rgb]{0.95,0.95,0.95}{
            \makecell[{{p{\linewidth}}}]{
                \texttt{\tiny{[P1$|$internal]}}
                \texttt{{[}attempt 1 – valid{]}} \\
\texttt{It looks like you're asking about a specific image or diagram, but I don't have any visual context or reference to what you're referring to. Could you please describe the image or provide more details about the layout (e.g., a map, a diagram, a room, a chart)? That way, I can help you answer your questions accurately.} \\
            }
        }
    }
    & & \\ \\

    \theutterance \stepcounter{utterance}  
    & \multicolumn{4}{p{0.6\linewidth}}{
        \cellcolor[rgb]{0.9,0.9,0.9}{
            \makecell[{{p{\linewidth}}}]{
                \texttt{\tiny{[P1$\rangle$GM]}}
                \texttt{It looks like you're asking about a specific image or diagram, but I don't have any visual context or reference to what you're referring to. Could you please describe the image or provide more details about the layout (e.g., a map, a diagram, a room, a chart)? That way, I can help you answer your questions accurately.} \\
            }
        }
    }
    & & \\ \\

    \theutterance \stepcounter{utterance}  
    & & \multicolumn{4}{p{0.6\linewidth}}{
        \cellcolor[rgb]{0.9,0.9,0.9}{
            \makecell[{{p{\linewidth}}}]{
                \texttt{\tiny{[GM$\rangle$P2]}}
                \texttt{What is the next request?} \\
            }
        }
    }
    & \\ \\

    \theutterance \stepcounter{utterance}  
    & & & \multicolumn{4}{p{0.6\linewidth}}{
        \cellcolor[rgb]{0.9,0.9,0.9}{
            \makecell[{{p{\linewidth}}}]{
                \texttt{\tiny{[GM$\langle$P2]}}
                \texttt{What is outside?} \\
            }
        }
    }
     \\ \\

    \theutterance \stepcounter{utterance}  
    & & \multicolumn{4}{p{0.6\linewidth}}{
        \cellcolor[rgb]{0.9,0.9,0.9}{
            \makecell[{{p{\linewidth}}}]{
                \texttt{\tiny{[P1$\langle$GM]}}
                \texttt{QUESTIONER: What is outside?} \\
            }
        }
    }
    & \\ \\

    \theutterance \stepcounter{utterance}  
    & \multicolumn{4}{p{0.45\linewidth}}{
        \cellcolor[rgb]{0.95,0.95,0.95}{
            \makecell[{{p{\linewidth}}}]{
                \texttt{\tiny{[P1$|$internal]}}
                \texttt{{[}attempt 1 – valid{]}} \\
\texttt{It looks like you're asking a series of questions that might be based on a visual or spatial context—like a map, diagram, or layout—but I don’t have any image or description to reference. Could you please provide more details or clarify what you're referring to? For example:} \\
\\ 
\texttt{{-} Are you asking about a specific map, room, diagram, or object?} \\
\texttt{{-} Is this related to a game, puzzle, or real{-}world location?} \\
\\ 
\texttt{Once I know the context, I’ll be happy to help answer your questions!} \\
            }
        }
    }
    & & \\ \\

    \theutterance \stepcounter{utterance}  
    & \multicolumn{4}{p{0.6\linewidth}}{
        \cellcolor[rgb]{0.9,0.9,0.9}{
            \makecell[{{p{\linewidth}}}]{
                \texttt{\tiny{[P1$\rangle$GM]}}
                \texttt{It looks like you're asking a series of questions that might be based on a visual or spatial context—like a map, diagram, or layout—but I don’t have any image or description to reference. Could you please provide more details or clarify what you're referring to? For example:} \\
\\ 
\texttt{{-} Are you asking about a specific map, room, diagram, or object?} \\
\texttt{{-} Is this related to a game, puzzle, or real{-}world location?} \\
\\ 
\texttt{Once I know the context, I’ll be happy to help answer your questions!} \\
            }
        }
    }
    & & \\ \\

    \theutterance \stepcounter{utterance}  
    & & \multicolumn{4}{p{0.6\linewidth}}{
        \cellcolor[rgb]{0.9,0.9,0.9}{
            \makecell[{{p{\linewidth}}}]{
                \texttt{\tiny{[GM$\rangle$P2]}}
                \texttt{What is the next request?} \\
            }
        }
    }
    & \\ \\

    \theutterance \stepcounter{utterance}  
    & & & \multicolumn{4}{p{0.6\linewidth}}{
        \cellcolor[rgb]{0.9,0.9,0.9}{
            \makecell[{{p{\linewidth}}}]{
                \texttt{\tiny{[GM$\langle$P2]}}
                \texttt{What is everywhere?} \\
            }
        }
    }
     \\ \\

    \theutterance \stepcounter{utterance}  
    & & \multicolumn{4}{p{0.6\linewidth}}{
        \cellcolor[rgb]{0.9,0.9,0.9}{
            \makecell[{{p{\linewidth}}}]{
                \texttt{\tiny{[P1$\langle$GM]}}
                \texttt{QUESTIONER: What is everywhere?} \\
            }
        }
    }
    & \\ \\

    \theutterance \stepcounter{utterance}  
    & \multicolumn{4}{p{0.45\linewidth}}{
        \cellcolor[rgb]{0.95,0.95,0.95}{
            \makecell[{{p{\linewidth}}}]{
                \texttt{\tiny{[P1$|$internal]}}
                \texttt{{[}attempt 1 – valid{]}} \\
\texttt{It seems like you're asking a series of questions that might be based on a specific context—perhaps a diagram, a map, a room, or a metaphorical scenario. Since I don’t have any visual or contextual information, I’ll offer a general, thoughtful interpretation of each question, assuming a symbolic or spatial framework (like a compass, a house, or a philosophical model):} \\
\\ 
\texttt{1. **What is at the top?**  } \\
\texttt{   *Answer:* The sky, the heavens, or the highest point—symbolizing aspiration, authority, or the divine.} \\
\\ 
\texttt{2. **What is in the northwest?**  } \\
\texttt{   *Answer:* A direction often associated with mystery, introspection, or the unknown—perhaps the land of the unseen or the realm of dreams.} \\
\\ 
\texttt{3. **What is at the center?**  } \\
\texttt{   *Answer:* The heart, the core, or the source—where everything converges and energy radiates from.} \\
\\ 
\texttt{4. **What is in the southwest?**  } \\
\texttt{   *Answer:* The place of grounding, stability, and earth—where roots grow and traditions are nurtured.} \\
\\ 
\texttt{5. **What is on your left?**  } \\
\texttt{   *Answer:* The side of the past, the intuitive, or the feminine—often associated with the subconscious or the unknown.} \\
\\ 
\texttt{6. **What is outside?**  } \\
\texttt{   *Answer:* The world beyond, the unknown, freedom, or the vastness of possibility.} \\
\\ 
\texttt{7. **What is everywhere?**  } \\
\texttt{   *} \\
            }
        }
    }
    & & \\ \\

    \theutterance \stepcounter{utterance}  
    & \multicolumn{4}{p{0.6\linewidth}}{
        \cellcolor[rgb]{0.9,0.9,0.9}{
            \makecell[{{p{\linewidth}}}]{
                \texttt{\tiny{[P1$\rangle$GM]}}
                \texttt{It seems like you're asking a series of questions that might be based on a specific context—perhaps a diagram, a map, a room, or a metaphorical scenario. Since I don’t have any visual or contextual information, I’ll offer a general, thoughtful interpretation of each question, assuming a symbolic or spatial framework (like a compass, a house, or a philosophical model):} \\
\\ 
\texttt{1. **What is at the top?**  } \\
\texttt{   *Answer:* The sky, the heavens, or the highest point—symbolizing aspiration, authority, or the divine.} \\
\\ 
\texttt{2. **What is in the northwest?**  } \\
\texttt{   *Answer:* A direction often associated with mystery, introspection, or the unknown—perhaps the land of the unseen or the realm of dreams.} \\
\\ 
\texttt{3. **What is at the center?**  } \\
\texttt{   *Answer:* The heart, the core, or the source—where everything converges and energy radiates from.} \\
\\ 
\texttt{4. **What is in the southwest?**  } \\
\texttt{   *Answer:* The place of grounding, stability, and earth—where roots grow and traditions are nurtured.} \\
\\ 
\texttt{5. **What is on your left?**  } \\
\texttt{   *Answer:* The side of the past, the intuitive, or the feminine—often associated with the subconscious or the unknown.} \\
\\ 
\texttt{6. **What is outside?**  } \\
\texttt{   *Answer:* The world beyond, the unknown, freedom, or the vastness of possibility.} \\
\\ 
\texttt{7. **What is everywhere?**  } \\
\texttt{   *} \\
            }
        }
    }
    & & \\ \\

    \theutterance \stepcounter{utterance}  
    & & \multicolumn{4}{p{0.6\linewidth}}{
        \cellcolor[rgb]{0.9,0.9,0.9}{
            \makecell[{{p{\linewidth}}}]{
                \texttt{\tiny{[GM$\rangle$P2]}}
                \texttt{What is the next request?} \\
            }
        }
    }
    & \\ \\

    \theutterance \stepcounter{utterance}  
    & & & \multicolumn{4}{p{0.6\linewidth}}{
        \cellcolor[rgb]{0.9,0.9,0.9}{
            \makecell[{{p{\linewidth}}}]{
                \texttt{\tiny{[GM$\langle$P2]}}
                \texttt{What is on your right?} \\
            }
        }
    }
     \\ \\

    \theutterance \stepcounter{utterance}  
    & & \multicolumn{4}{p{0.6\linewidth}}{
        \cellcolor[rgb]{0.9,0.9,0.9}{
            \makecell[{{p{\linewidth}}}]{
                \texttt{\tiny{[P1$\langle$GM]}}
                \texttt{QUESTIONER: What is on your right?} \\
            }
        }
    }
    & \\ \\

    \theutterance \stepcounter{utterance}  
    & \multicolumn{4}{p{0.45\linewidth}}{
        \cellcolor[rgb]{0.95,0.95,0.95}{
            \makecell[{{p{\linewidth}}}]{
                \texttt{\tiny{[P1$|$internal]}}
                \texttt{{[}attempt 1 – valid{]}} \\
\texttt{It seems like you're asking a series of spatial or philosophical questions, possibly in the context of a riddle, meditation, or symbolic framework. Since there's no visual or contextual reference provided (like a map, diagram, or environment), I'll interpret these questions in a general, metaphorical, or universal way—common in spiritual or philosophical traditions.} \\
\\ 
\texttt{Here’s a possible interpretation:} \\
\\ 
\texttt{1. **What is at the top?**  } \\
\texttt{   *The sky* — or *the heavens* — symbolizing the divine, the infinite, or the highest point.} \\
\\ 
\texttt{2. **What is in the northwest?**  } \\
\texttt{   *The unknown* — or *the direction of mystery* — often associated with introspection, the past, or hidden knowledge in some traditions.} \\
\\ 
\texttt{3. **What is at the center?**  } \\
\texttt{   *The self* — or *the heart* — representing balance, presence, and the core of being.} \\
\\ 
\texttt{4. **What is in the southwest?**  } \\
\texttt{   *The earth* — or *the foundation* — symbolizing stability, grounding, and the material world.} \\
\\ 
\texttt{5. **What is on your left?**  } \\
\texttt{   *The past* — or *the feminine* — often associated with intuition, memory, and the subconscious.} \\
\\ 
\texttt{6. **What is outside?**  } \\
\texttt{   *The world* — or *the unknown* — representing the external environment, other people, and the broader reality beyond the self.} \\
\\ 
\texttt{7. **What is everywhere?} \\
            }
        }
    }
    & & \\ \\

    \theutterance \stepcounter{utterance}  
    & \multicolumn{4}{p{0.6\linewidth}}{
        \cellcolor[rgb]{0.9,0.9,0.9}{
            \makecell[{{p{\linewidth}}}]{
                \texttt{\tiny{[P1$\rangle$GM]}}
                \texttt{It seems like you're asking a series of spatial or philosophical questions, possibly in the context of a riddle, meditation, or symbolic framework. Since there's no visual or contextual reference provided (like a map, diagram, or environment), I'll interpret these questions in a general, metaphorical, or universal way—common in spiritual or philosophical traditions.} \\
\\ 
\texttt{Here’s a possible interpretation:} \\
\\ 
\texttt{1. **What is at the top?**  } \\
\texttt{   *The sky* — or *the heavens* — symbolizing the divine, the infinite, or the highest point.} \\
\\ 
\texttt{2. **What is in the northwest?**  } \\
\texttt{   *The unknown* — or *the direction of mystery* — often associated with introspection, the past, or hidden knowledge in some traditions.} \\
\\ 
\texttt{3. **What is at the center?**  } \\
\texttt{   *The self* — or *the heart* — representing balance, presence, and the core of being.} \\
\\ 
\texttt{4. **What is in the southwest?**  } \\
\texttt{   *The earth* — or *the foundation* — symbolizing stability, grounding, and the material world.} \\
\\ 
\texttt{5. **What is on your left?**  } \\
\texttt{   *The past* — or *the feminine* — often associated with intuition, memory, and the subconscious.} \\
\\ 
\texttt{6. **What is outside?**  } \\
\texttt{   *The world* — or *the unknown* — representing the external environment, other people, and the broader reality beyond the self.} \\
\\ 
\texttt{7. **What is everywhere?} \\
            }
        }
    }
    & & \\ \\

    \theutterance \stepcounter{utterance}  
    & & \multicolumn{4}{p{0.6\linewidth}}{
        \cellcolor[rgb]{0.9,0.9,0.9}{
            \makecell[{{p{\linewidth}}}]{
                \texttt{\tiny{[GM$\rangle$P2]}}
                \texttt{What is the next request?} \\
            }
        }
    }
    & \\ \\

    \theutterance \stepcounter{utterance}  
    & & & \multicolumn{4}{p{0.6\linewidth}}{
        \cellcolor[rgb]{0.9,0.9,0.9}{
            \makecell[{{p{\linewidth}}}]{
                \texttt{\tiny{[GM$\langle$P2]}}
                \texttt{What is nowhere?} \\
            }
        }
    }
     \\ \\

    \theutterance \stepcounter{utterance}  
    & & \multicolumn{4}{p{0.6\linewidth}}{
        \cellcolor[rgb]{0.9,0.9,0.9}{
            \makecell[{{p{\linewidth}}}]{
                \texttt{\tiny{[P1$\langle$GM]}}
                \texttt{QUESTIONER: What is nowhere?} \\
            }
        }
    }
    & \\ \\

    \theutterance \stepcounter{utterance}  
    & \multicolumn{4}{p{0.45\linewidth}}{
        \cellcolor[rgb]{0.95,0.95,0.95}{
            \makecell[{{p{\linewidth}}}]{
                \texttt{\tiny{[P1$|$internal]}}
                \texttt{{[}attempt 1 – valid{]}} \\
\texttt{These questions seem to be philosophical or metaphorical in nature, possibly inspired by riddles or meditative inquiries. Let's explore each one with a thoughtful, interpretive lens:} \\
\\ 
\texttt{1. **What is at the top?**  } \\
\texttt{   *Answer:* The sky, the heavens, or perhaps the source of all things—symbolizing the highest point, the divine, or the ultimate truth.} \\
\\ 
\texttt{2. **What is in the northwest?**  } \\
\texttt{   *Answer:* A direction, a place, or a mystery—perhaps the unknown, the hidden, or the realm of introspection. In some traditions, the northwest is associated with the element of water or the spirit of the unknown.} \\
\\ 
\texttt{3. **What is at the center?**  } \\
\texttt{   *Answer:* The heart, the self, the core of being—where everything converges and from which all things radiate.} \\
\\ 
\texttt{4. **What is in the southwest?**  } \\
\texttt{   *Answer:* The place of grounding, the earth, or the realm of the subconscious. In some systems, it's associated with the element of earth or the past.} \\
\\ 
\texttt{5. **What is on your left?**  } \\
\texttt{   *Answer:* The past, the intuitive, the feminine, or the unseen—often symbolizing what is behind or to the side, not directly in front.} \\
\\ 
\texttt{6. **What is outside?**  } \\
\texttt{   *Answer:* The world, the other, the unknown, or the boundary beyond the self—what lies beyond our immediate experience} \\
            }
        }
    }
    & & \\ \\

    \theutterance \stepcounter{utterance}  
    & \multicolumn{4}{p{0.6\linewidth}}{
        \cellcolor[rgb]{0.9,0.9,0.9}{
            \makecell[{{p{\linewidth}}}]{
                \texttt{\tiny{[P1$\rangle$GM]}}
                \texttt{These questions seem to be philosophical or metaphorical in nature, possibly inspired by riddles or meditative inquiries. Let's explore each one with a thoughtful, interpretive lens:} \\
\\ 
\texttt{1. **What is at the top?**  } \\
\texttt{   *Answer:* The sky, the heavens, or perhaps the source of all things—symbolizing the highest point, the divine, or the ultimate truth.} \\
\\ 
\texttt{2. **What is in the northwest?**  } \\
\texttt{   *Answer:* A direction, a place, or a mystery—perhaps the unknown, the hidden, or the realm of introspection. In some traditions, the northwest is associated with the element of water or the spirit of the unknown.} \\
\\ 
\texttt{3. **What is at the center?**  } \\
\texttt{   *Answer:* The heart, the self, the core of being—where everything converges and from which all things radiate.} \\
\\ 
\texttt{4. **What is in the southwest?**  } \\
\texttt{   *Answer:* The place of grounding, the earth, or the realm of the subconscious. In some systems, it's associated with the element of earth or the past.} \\
\\ 
\texttt{5. **What is on your left?**  } \\
\texttt{   *Answer:* The past, the intuitive, the feminine, or the unseen—often symbolizing what is behind or to the side, not directly in front.} \\
\\ 
\texttt{6. **What is outside?**  } \\
\texttt{   *Answer:* The world, the other, the unknown, or the boundary beyond the self—what lies beyond our immediate experience} \\
            }
        }
    }
    & & \\ \\

    \theutterance \stepcounter{utterance}  
    & & \multicolumn{4}{p{0.6\linewidth}}{
        \cellcolor[rgb]{0.9,0.9,0.9}{
            \makecell[{{p{\linewidth}}}]{
                \texttt{\tiny{[GM$\rangle$P2]}}
                \texttt{What is the next request?} \\
            }
        }
    }
    & \\ \\

    \theutterance \stepcounter{utterance}  
    & & & \multicolumn{4}{p{0.6\linewidth}}{
        \cellcolor[rgb]{0.9,0.9,0.9}{
            \makecell[{{p{\linewidth}}}]{
                \texttt{\tiny{[GM$\langle$P2]}}
                \texttt{What is at the bottom?} \\
            }
        }
    }
     \\ \\

    \theutterance \stepcounter{utterance}  
    & & \multicolumn{4}{p{0.6\linewidth}}{
        \cellcolor[rgb]{0.9,0.9,0.9}{
            \makecell[{{p{\linewidth}}}]{
                \texttt{\tiny{[P1$\langle$GM]}}
                \texttt{QUESTIONER: What is at the bottom?} \\
            }
        }
    }
    & \\ \\

    \theutterance \stepcounter{utterance}  
    & \multicolumn{4}{p{0.45\linewidth}}{
        \cellcolor[rgb]{0.95,0.95,0.95}{
            \makecell[{{p{\linewidth}}}]{
                \texttt{\tiny{[P1$|$internal]}}
                \texttt{{[}attempt 1 – valid{]}} \\
\texttt{These questions appear to be a series of riddles or philosophical inquiries, possibly inspired by Zen koans or meditative puzzles. Let's explore each one with a thoughtful, interpretive lens—often the answers are not literal but symbolic or conceptual.} \\
\\ 
\texttt{{-}{-}{-}} \\
\\ 
\texttt{**1. What is at the top?**  } \\
\texttt{*Answer:* The sky, or the heavens.  } \\
\texttt{*Interpretation:* The highest point, the limit of perception, or the divine. In a metaphorical sense, it could be "nothing" — because the top of the universe is beyond reach, or "awareness" — the highest state of consciousness.} \\
\\ 
\texttt{{-}{-}{-}} \\
\\ 
\texttt{**2. What is in the northwest?**  } \\
\texttt{*Answer:* A direction.  } \\
\texttt{*Interpretation:* The northwest is a compass point, but in a deeper sense, it may represent the unknown, the hidden, or the place where the mind wanders. It could also be "the wind" — a symbol of change or movement.} \\
\\ 
\texttt{{-}{-}{-}} \\
\\ 
\texttt{**3. What is at the center?**  } \\
\texttt{*Answer:* The heart, the core, or the self.  } \\
\texttt{*Interpretation:* The center is where everything converges — the still point in the whirlwind, the "I" in the midst of experience. It could also be "now" — the present moment.} \\
\\ 
\texttt{{-}{-}{-}} \\
\\ 
\texttt{**4. What is in the southwest?**  } \\
\texttt{*Answer:* A direction.  } \\
\texttt{*Interpretation:* The southwest is often associated with earth, grounding, and} \\
            }
        }
    }
    & & \\ \\

    \theutterance \stepcounter{utterance}  
    & \multicolumn{4}{p{0.6\linewidth}}{
        \cellcolor[rgb]{0.9,0.9,0.9}{
            \makecell[{{p{\linewidth}}}]{
                \texttt{\tiny{[P1$\rangle$GM]}}
                \texttt{These questions appear to be a series of riddles or philosophical inquiries, possibly inspired by Zen koans or meditative puzzles. Let's explore each one with a thoughtful, interpretive lens—often the answers are not literal but symbolic or conceptual.} \\
\\ 
\texttt{{-}{-}{-}} \\
\\ 
\texttt{**1. What is at the top?**  } \\
\texttt{*Answer:* The sky, or the heavens.  } \\
\texttt{*Interpretation:* The highest point, the limit of perception, or the divine. In a metaphorical sense, it could be "nothing" — because the top of the universe is beyond reach, or "awareness" — the highest state of consciousness.} \\
\\ 
\texttt{{-}{-}{-}} \\
\\ 
\texttt{**2. What is in the northwest?**  } \\
\texttt{*Answer:* A direction.  } \\
\texttt{*Interpretation:* The northwest is a compass point, but in a deeper sense, it may represent the unknown, the hidden, or the place where the mind wanders. It could also be "the wind" — a symbol of change or movement.} \\
\\ 
\texttt{{-}{-}{-}} \\
\\ 
\texttt{**3. What is at the center?**  } \\
\texttt{*Answer:* The heart, the core, or the self.  } \\
\texttt{*Interpretation:* The center is where everything converges — the still point in the whirlwind, the "I" in the midst of experience. It could also be "now" — the present moment.} \\
\\ 
\texttt{{-}{-}{-}} \\
\\ 
\texttt{**4. What is in the southwest?**  } \\
\texttt{*Answer:* A direction.  } \\
\texttt{*Interpretation:* The southwest is often associated with earth, grounding, and} \\
            }
        }
    }
    & & \\ \\

    \theutterance \stepcounter{utterance}  
    & & \multicolumn{4}{p{0.6\linewidth}}{
        \cellcolor[rgb]{0.9,0.9,0.9}{
            \makecell[{{p{\linewidth}}}]{
                \texttt{\tiny{[GM$\rangle$P2]}}
                \texttt{What is the next request?} \\
            }
        }
    }
    & \\ \\

    \theutterance \stepcounter{utterance}  
    & & & \multicolumn{4}{p{0.6\linewidth}}{
        \cellcolor[rgb]{0.9,0.9,0.9}{
            \makecell[{{p{\linewidth}}}]{
                \texttt{\tiny{[GM$\langle$P2]}}
                \texttt{What is in the southeast?} \\
            }
        }
    }
     \\ \\

    \theutterance \stepcounter{utterance}  
    & & \multicolumn{4}{p{0.6\linewidth}}{
        \cellcolor[rgb]{0.9,0.9,0.9}{
            \makecell[{{p{\linewidth}}}]{
                \texttt{\tiny{[P1$\langle$GM]}}
                \texttt{QUESTIONER: What is in the southeast?} \\
            }
        }
    }
    & \\ \\

    \theutterance \stepcounter{utterance}  
    & \multicolumn{4}{p{0.45\linewidth}}{
        \cellcolor[rgb]{0.95,0.95,0.95}{
            \makecell[{{p{\linewidth}}}]{
                \texttt{\tiny{[P1$|$internal]}}
                \texttt{{[}attempt 1 – valid{]}} \\
\texttt{These questions appear to be a series of riddles or philosophical inquiries, possibly inspired by Zen koans or meditative puzzles. They are designed to provoke thought and challenge conventional thinking. Let's explore each one with a reflective, interpretive lens:} \\
\\ 
\texttt{{-}{-}{-}} \\
\\ 
\texttt{**1. What is at the top?**  } \\
\texttt{*Answer:* The sky, the heavens, or perhaps "nothing" — since the top is often seen as the limit, beyond which there is no more space. In a metaphorical sense, it could be "infinity" or "the unknown."} \\
\\ 
\texttt{**2. What is in the northwest?**  } \\
\texttt{*Answer:* A direction — a point on the compass. But in a deeper sense, it might be "the unknown," "the hidden," or "the place where the wind comes from." Or, in a poetic sense, "the land of the forgotten."} \\
\\ 
\texttt{**3. What is at the center?**  } \\
\texttt{*Answer:* The heart, the core, the self, or "the present moment." In a physical sense, it could be "the eye of the storm" — calm amidst chaos.} \\
\\ 
\texttt{**4. What is in the southwest?**  } \\
\texttt{*Answer:* A direction — perhaps "the land of the setting sun," "the place of endings," or "the realm of the subconscious." In some traditions, it's associated with the element of water or the emotion of sorrow.} \\
\\ 
\texttt{**5. What is on your left?**  } \\
\texttt{*Answer:* The side opposite your} \\
            }
        }
    }
    & & \\ \\

    \theutterance \stepcounter{utterance}  
    & \multicolumn{4}{p{0.6\linewidth}}{
        \cellcolor[rgb]{0.9,0.9,0.9}{
            \makecell[{{p{\linewidth}}}]{
                \texttt{\tiny{[P1$\rangle$GM]}}
                \texttt{These questions appear to be a series of riddles or philosophical inquiries, possibly inspired by Zen koans or meditative puzzles. They are designed to provoke thought and challenge conventional thinking. Let's explore each one with a reflective, interpretive lens:} \\
\\ 
\texttt{{-}{-}{-}} \\
\\ 
\texttt{**1. What is at the top?**  } \\
\texttt{*Answer:* The sky, the heavens, or perhaps "nothing" — since the top is often seen as the limit, beyond which there is no more space. In a metaphorical sense, it could be "infinity" or "the unknown."} \\
\\ 
\texttt{**2. What is in the northwest?**  } \\
\texttt{*Answer:* A direction — a point on the compass. But in a deeper sense, it might be "the unknown," "the hidden," or "the place where the wind comes from." Or, in a poetic sense, "the land of the forgotten."} \\
\\ 
\texttt{**3. What is at the center?**  } \\
\texttt{*Answer:* The heart, the core, the self, or "the present moment." In a physical sense, it could be "the eye of the storm" — calm amidst chaos.} \\
\\ 
\texttt{**4. What is in the southwest?**  } \\
\texttt{*Answer:* A direction — perhaps "the land of the setting sun," "the place of endings," or "the realm of the subconscious." In some traditions, it's associated with the element of water or the emotion of sorrow.} \\
\\ 
\texttt{**5. What is on your left?**  } \\
\texttt{*Answer:* The side opposite your} \\
            }
        }
    }
    & & \\ \\

    \theutterance \stepcounter{utterance}  
    & & \multicolumn{4}{p{0.6\linewidth}}{
        \cellcolor[rgb]{0.9,0.9,0.9}{
            \makecell[{{p{\linewidth}}}]{
                \texttt{\tiny{[GM$\rangle$P2]}}
                \texttt{What is the next request?} \\
            }
        }
    }
    & \\ \\

    \theutterance \stepcounter{utterance}  
    & & & \multicolumn{4}{p{0.6\linewidth}}{
        \cellcolor[rgb]{0.9,0.9,0.9}{
            \makecell[{{p{\linewidth}}}]{
                \texttt{\tiny{[GM$\langle$P2]}}
                \texttt{What is there?} \\
            }
        }
    }
     \\ \\

    \theutterance \stepcounter{utterance}  
    & & \multicolumn{4}{p{0.6\linewidth}}{
        \cellcolor[rgb]{0.9,0.9,0.9}{
            \makecell[{{p{\linewidth}}}]{
                \texttt{\tiny{[P1$\langle$GM]}}
                \texttt{QUESTIONER: What is there?} \\
            }
        }
    }
    & \\ \\

    \theutterance \stepcounter{utterance}  
    & \multicolumn{4}{p{0.45\linewidth}}{
        \cellcolor[rgb]{0.95,0.95,0.95}{
            \makecell[{{p{\linewidth}}}]{
                \texttt{\tiny{[P1$|$internal]}}
                \texttt{{[}attempt 1 – valid{]}} \\
\texttt{These questions seem to be a series of riddles or philosophical inquiries, possibly inspired by a spatial or metaphysical framework. Let's interpret them in a poetic, symbolic, or philosophical way, as they don't refer to a specific image or context.} \\
\\ 
\texttt{Here’s a possible interpretation:} \\
\\ 
\texttt{1. **What is at the top?**  } \\
\texttt{   *The sky* — or *the heavens*, representing the highest point, the divine, or the unknown.} \\
\\ 
\texttt{2. **What is in the northwest?**  } \\
\texttt{   *The unknown* — or *the cold, the distant*, symbolizing mystery, isolation, or the unexplored.} \\
\\ 
\texttt{3. **What is at the center?**  } \\
\texttt{   *The self* — or *the heart*, representing the core of being, the point of balance.} \\
\\ 
\texttt{4. **What is in the southwest?**  } \\
\texttt{   *The earth* — or *the grounded*, symbolizing stability, nature, or the subconscious.} \\
\\ 
\texttt{5. **What is on your left?**  } \\
\texttt{   *The past* — or *memory*, representing what has been, what is behind us.} \\
\\ 
\texttt{6. **What is outside?**  } \\
\texttt{   *The world* — or *the unknown*, the space beyond our immediate experience.} \\
\\ 
\texttt{7. **What is everywhere?**  } \\
\texttt{   *The present* — or *awareness*, the space where all things exist simultaneously.} \\
\\ 
\texttt{8. **What is on your right?**  } \\
\texttt{   *The future* — or *possibility*,} \\
            }
        }
    }
    & & \\ \\

    \theutterance \stepcounter{utterance}  
    & \multicolumn{4}{p{0.6\linewidth}}{
        \cellcolor[rgb]{0.9,0.9,0.9}{
            \makecell[{{p{\linewidth}}}]{
                \texttt{\tiny{[P1$\rangle$GM]}}
                \texttt{These questions seem to be a series of riddles or philosophical inquiries, possibly inspired by a spatial or metaphysical framework. Let's interpret them in a poetic, symbolic, or philosophical way, as they don't refer to a specific image or context.} \\
\\ 
\texttt{Here’s a possible interpretation:} \\
\\ 
\texttt{1. **What is at the top?**  } \\
\texttt{   *The sky* — or *the heavens*, representing the highest point, the divine, or the unknown.} \\
\\ 
\texttt{2. **What is in the northwest?**  } \\
\texttt{   *The unknown* — or *the cold, the distant*, symbolizing mystery, isolation, or the unexplored.} \\
\\ 
\texttt{3. **What is at the center?**  } \\
\texttt{   *The self* — or *the heart*, representing the core of being, the point of balance.} \\
\\ 
\texttt{4. **What is in the southwest?**  } \\
\texttt{   *The earth* — or *the grounded*, symbolizing stability, nature, or the subconscious.} \\
\\ 
\texttt{5. **What is on your left?**  } \\
\texttt{   *The past* — or *memory*, representing what has been, what is behind us.} \\
\\ 
\texttt{6. **What is outside?**  } \\
\texttt{   *The world* — or *the unknown*, the space beyond our immediate experience.} \\
\\ 
\texttt{7. **What is everywhere?**  } \\
\texttt{   *The present* — or *awareness*, the space where all things exist simultaneously.} \\
\\ 
\texttt{8. **What is on your right?**  } \\
\texttt{   *The future* — or *possibility*,} \\
            }
        }
    }
    & & \\ \\

    \theutterance \stepcounter{utterance}  
    & & \multicolumn{4}{p{0.6\linewidth}}{
        \cellcolor[rgb]{0.9,0.9,0.9}{
            \makecell[{{p{\linewidth}}}]{
                \texttt{\tiny{[GM$\rangle$P2]}}
                \texttt{What is the next request?} \\
            }
        }
    }
    & \\ \\

    \theutterance \stepcounter{utterance}  
    & & & \multicolumn{4}{p{0.6\linewidth}}{
        \cellcolor[rgb]{0.9,0.9,0.9}{
            \makecell[{{p{\linewidth}}}]{
                \texttt{\tiny{[GM$\langle$P2]}}
                \texttt{What is inside?} \\
            }
        }
    }
     \\ \\

    \theutterance \stepcounter{utterance}  
    & & \multicolumn{4}{p{0.6\linewidth}}{
        \cellcolor[rgb]{0.9,0.9,0.9}{
            \makecell[{{p{\linewidth}}}]{
                \texttt{\tiny{[P1$\langle$GM]}}
                \texttt{QUESTIONER: What is inside?} \\
            }
        }
    }
    & \\ \\

    \theutterance \stepcounter{utterance}  
    & \multicolumn{4}{p{0.45\linewidth}}{
        \cellcolor[rgb]{0.95,0.95,0.95}{
            \makecell[{{p{\linewidth}}}]{
                \texttt{\tiny{[P1$|$internal]}}
                \texttt{{[}attempt 1 – valid{]}} \\
\texttt{These questions appear to be a series of riddles or philosophical inquiries, possibly inspired by Zen koans or meditative puzzles. They are designed to provoke thought and challenge conventional thinking. Let's explore each one with a reflective, open{-}ended interpretation:} \\
\\ 
\texttt{{-}{-}{-}} \\
\\ 
\texttt{**1. What is at the top?**  } \\
\texttt{*Answer:* The sky, the heavens, or perhaps "nothing" — since the top is often seen as the limit, beyond which there is no more.  } \\
\texttt{*Reflection:* The top may be where we look for answers, but it's also where we find silence.} \\
\\ 
\texttt{{-}{-}{-}} \\
\\ 
\texttt{**2. What is in the northwest?**  } \\
\texttt{*Answer:* A direction — a point on the compass. But perhaps it's "the unknown" or "the place where the wind comes from."  } \\
\texttt{*Reflection:* The northwest is a place of mystery, often associated with cold, change, or the edge of the known world.} \\
\\ 
\texttt{{-}{-}{-}} \\
\\ 
\texttt{**3. What is at the center?**  } \\
\texttt{*Answer:* The heart, the core, the self, or "the still point of the turning world."  } \\
\texttt{*Reflection:* The center is where everything converges — the place of balance and awareness.} \\
\\ 
\texttt{{-}{-}{-}} \\
\\ 
\texttt{**4. What is in the southwest?**  } \\
\texttt{*Answer:* The warmth of the sun, the land of shadows, or "the place where the earth meets the sky."  } \\
\texttt{*Reflection:* The southwest is often associated with grounding, intuition, and the feminine.} \\
\\ 
\texttt{{-}{-}{-}} \\
\\ 
\texttt{**5. What is on} \\
            }
        }
    }
    & & \\ \\

    \theutterance \stepcounter{utterance}  
    & \multicolumn{4}{p{0.6\linewidth}}{
        \cellcolor[rgb]{0.9,0.9,0.9}{
            \makecell[{{p{\linewidth}}}]{
                \texttt{\tiny{[P1$\rangle$GM]}}
                \texttt{These questions appear to be a series of riddles or philosophical inquiries, possibly inspired by Zen koans or meditative puzzles. They are designed to provoke thought and challenge conventional thinking. Let's explore each one with a reflective, open{-}ended interpretation:} \\
\\ 
\texttt{{-}{-}{-}} \\
\\ 
\texttt{**1. What is at the top?**  } \\
\texttt{*Answer:* The sky, the heavens, or perhaps "nothing" — since the top is often seen as the limit, beyond which there is no more.  } \\
\texttt{*Reflection:* The top may be where we look for answers, but it's also where we find silence.} \\
\\ 
\texttt{{-}{-}{-}} \\
\\ 
\texttt{**2. What is in the northwest?**  } \\
\texttt{*Answer:* A direction — a point on the compass. But perhaps it's "the unknown" or "the place where the wind comes from."  } \\
\texttt{*Reflection:* The northwest is a place of mystery, often associated with cold, change, or the edge of the known world.} \\
\\ 
\texttt{{-}{-}{-}} \\
\\ 
\texttt{**3. What is at the center?**  } \\
\texttt{*Answer:* The heart, the core, the self, or "the still point of the turning world."  } \\
\texttt{*Reflection:* The center is where everything converges — the place of balance and awareness.} \\
\\ 
\texttt{{-}{-}{-}} \\
\\ 
\texttt{**4. What is in the southwest?**  } \\
\texttt{*Answer:* The warmth of the sun, the land of shadows, or "the place where the earth meets the sky."  } \\
\texttt{*Reflection:* The southwest is often associated with grounding, intuition, and the feminine.} \\
\\ 
\texttt{{-}{-}{-}} \\
\\ 
\texttt{**5. What is on} \\
            }
        }
    }
    & & \\ \\

    \theutterance \stepcounter{utterance}  
    & & \multicolumn{4}{p{0.6\linewidth}}{
        \cellcolor[rgb]{0.9,0.9,0.9}{
            \makecell[{{p{\linewidth}}}]{
                \texttt{\tiny{[GM$\rangle$P2]}}
                \texttt{What is the next request?} \\
            }
        }
    }
    & \\ \\

    \theutterance \stepcounter{utterance}  
    & & & \multicolumn{4}{p{0.6\linewidth}}{
        \cellcolor[rgb]{0.9,0.9,0.9}{
            \makecell[{{p{\linewidth}}}]{
                \texttt{\tiny{[GM$\langle$P2]}}
                \texttt{What is in the northeast?} \\
            }
        }
    }
     \\ \\

    \theutterance \stepcounter{utterance}  
    & & \multicolumn{4}{p{0.6\linewidth}}{
        \cellcolor[rgb]{0.9,0.9,0.9}{
            \makecell[{{p{\linewidth}}}]{
                \texttt{\tiny{[P1$\langle$GM]}}
                \texttt{QUESTIONER: What is in the northeast?} \\
            }
        }
    }
    & \\ \\

    \theutterance \stepcounter{utterance}  
    & \multicolumn{4}{p{0.45\linewidth}}{
        \cellcolor[rgb]{0.95,0.95,0.95}{
            \makecell[{{p{\linewidth}}}]{
                \texttt{\tiny{[P1$|$internal]}}
                \texttt{{[}attempt 1 – valid{]}} \\
\texttt{These questions appear to be a series of riddles or philosophical inquiries, possibly inspired by Zen koans or meditative puzzles. They are structured to provoke thought about perception, reality, and the nature of existence. Let's explore each one with a reflective, interpretive lens—since literal answers may not be the point.} \\
\\ 
\texttt{{-}{-}{-}} \\
\\ 
\texttt{**1. What is at the top?**  } \\
\texttt{*Answer:* The sky, the heavens, or perhaps "nothing"—since the top is often where things end, or where the mind reaches for the infinite.} \\
\\ 
\texttt{**2. What is in the northwest?**  } \\
\texttt{*Answer:* A direction, a point on the compass—perhaps "the unknown" or "the place where the wind comes from." Or, in a metaphorical sense, "the past" or "the hidden."} \\
\\ 
\texttt{**3. What is at the center?**  } \\
\texttt{*Answer:* The heart, the core, the self, or "the present moment." In a physical sense, it could be "you" or "the observer."} \\
\\ 
\texttt{**4. What is in the southwest?**  } \\
\texttt{*Answer:* The place of the setting sun, the realm of endings, or "the land of shadows." In some traditions, it's associated with the element of earth or the color red.} \\
\\ 
\texttt{**5. What is on your left?**  } \\
\texttt{*Answer:* The side opposite your dominant hand—perhaps "the past" or "the intuitive." Or simply, "the other."} \\
\\ 
\texttt{**6. What} \\
            }
        }
    }
    & & \\ \\

    \theutterance \stepcounter{utterance}  
    & \multicolumn{4}{p{0.6\linewidth}}{
        \cellcolor[rgb]{0.9,0.9,0.9}{
            \makecell[{{p{\linewidth}}}]{
                \texttt{\tiny{[P1$\rangle$GM]}}
                \texttt{These questions appear to be a series of riddles or philosophical inquiries, possibly inspired by Zen koans or meditative puzzles. They are structured to provoke thought about perception, reality, and the nature of existence. Let's explore each one with a reflective, interpretive lens—since literal answers may not be the point.} \\
\\ 
\texttt{{-}{-}{-}} \\
\\ 
\texttt{**1. What is at the top?**  } \\
\texttt{*Answer:* The sky, the heavens, or perhaps "nothing"—since the top is often where things end, or where the mind reaches for the infinite.} \\
\\ 
\texttt{**2. What is in the northwest?**  } \\
\texttt{*Answer:* A direction, a point on the compass—perhaps "the unknown" or "the place where the wind comes from." Or, in a metaphorical sense, "the past" or "the hidden."} \\
\\ 
\texttt{**3. What is at the center?**  } \\
\texttt{*Answer:* The heart, the core, the self, or "the present moment." In a physical sense, it could be "you" or "the observer."} \\
\\ 
\texttt{**4. What is in the southwest?**  } \\
\texttt{*Answer:* The place of the setting sun, the realm of endings, or "the land of shadows." In some traditions, it's associated with the element of earth or the color red.} \\
\\ 
\texttt{**5. What is on your left?**  } \\
\texttt{*Answer:* The side opposite your dominant hand—perhaps "the past" or "the intuitive." Or simply, "the other."} \\
\\ 
\texttt{**6. What} \\
            }
        }
    }
    & & \\ \\

    \theutterance \stepcounter{utterance}  
    & & \multicolumn{4}{p{0.6\linewidth}}{
        \cellcolor[rgb]{0.9,0.9,0.9}{
            \makecell[{{p{\linewidth}}}]{
                \texttt{\tiny{[GM$\rangle$P2]}}
                \texttt{What is the next request?} \\
            }
        }
    }
    & \\ \\

    \theutterance \stepcounter{utterance}  
    & & & \multicolumn{4}{p{0.6\linewidth}}{
        \cellcolor[rgb]{0.9,0.9,0.9}{
            \makecell[{{p{\linewidth}}}]{
                \texttt{\tiny{[GM$\langle$P2]}}
                \texttt{What is here?} \\
            }
        }
    }
     \\ \\

    \theutterance \stepcounter{utterance}  
    & & \multicolumn{4}{p{0.6\linewidth}}{
        \cellcolor[rgb]{0.9,0.9,0.9}{
            \makecell[{{p{\linewidth}}}]{
                \texttt{\tiny{[P1$\langle$GM]}}
                \texttt{QUESTIONER: What is here?} \\
            }
        }
    }
    & \\ \\

    \theutterance \stepcounter{utterance}  
    & \multicolumn{4}{p{0.45\linewidth}}{
        \cellcolor[rgb]{0.95,0.95,0.95}{
            \makecell[{{p{\linewidth}}}]{
                \texttt{\tiny{[P1$|$internal]}}
                \texttt{{[}attempt 1 – valid{]}} \\
\texttt{These questions appear to be a series of riddles or philosophical inquiries, possibly inspired by Zen koans or meditative prompts. They are designed to provoke thought and reflection rather than to elicit literal answers. Let's explore them in a contemplative way:} \\
\\ 
\texttt{{-}{-}{-}} \\
\\ 
\texttt{**1. What is at the top?**  } \\
\texttt{*Answer:* The sky, the heavens, or perhaps "nothing"—since the top is often where things end, and beyond it lies the infinite.} \\
\\ 
\texttt{**2. What is in the northwest?**  } \\
\texttt{*Answer:* A direction, a point on the compass—perhaps "the unknown" or "the place where the wind comes from."} \\
\\ 
\texttt{**3. What is at the center?**  } \\
\texttt{*Answer:* The heart, the core, the self, or "the present moment."} \\
\\ 
\texttt{**4. What is in the southwest?**  } \\
\texttt{*Answer:* The place of the setting sun, the realm of endings, or "the quiet."} \\
\\ 
\texttt{**5. What is on your left?**  } \\
\texttt{*Answer:* The side of the body, or "the past" (if you're facing forward).} \\
\\ 
\texttt{**6. What is outside?**  } \\
\texttt{*Answer:* The world, the environment, or "the other."} \\
\\ 
\texttt{**7. What is everywhere?**  } \\
\texttt{*Answer:* Space, awareness, or "the mind."} \\
\\ 
\texttt{**8. What is on your right?**  } \\
\texttt{*Answer:* The side of the body, or "the future."} \\
\\ 
\texttt{**9. What is nowhere?**} \\
            }
        }
    }
    & & \\ \\

    \theutterance \stepcounter{utterance}  
    & \multicolumn{4}{p{0.6\linewidth}}{
        \cellcolor[rgb]{0.9,0.9,0.9}{
            \makecell[{{p{\linewidth}}}]{
                \texttt{\tiny{[P1$\rangle$GM]}}
                \texttt{These questions appear to be a series of riddles or philosophical inquiries, possibly inspired by Zen koans or meditative prompts. They are designed to provoke thought and reflection rather than to elicit literal answers. Let's explore them in a contemplative way:} \\
\\ 
\texttt{{-}{-}{-}} \\
\\ 
\texttt{**1. What is at the top?**  } \\
\texttt{*Answer:* The sky, the heavens, or perhaps "nothing"—since the top is often where things end, and beyond it lies the infinite.} \\
\\ 
\texttt{**2. What is in the northwest?**  } \\
\texttt{*Answer:* A direction, a point on the compass—perhaps "the unknown" or "the place where the wind comes from."} \\
\\ 
\texttt{**3. What is at the center?**  } \\
\texttt{*Answer:* The heart, the core, the self, or "the present moment."} \\
\\ 
\texttt{**4. What is in the southwest?**  } \\
\texttt{*Answer:* The place of the setting sun, the realm of endings, or "the quiet."} \\
\\ 
\texttt{**5. What is on your left?**  } \\
\texttt{*Answer:* The side of the body, or "the past" (if you're facing forward).} \\
\\ 
\texttt{**6. What is outside?**  } \\
\texttt{*Answer:* The world, the environment, or "the other."} \\
\\ 
\texttt{**7. What is everywhere?**  } \\
\texttt{*Answer:* Space, awareness, or "the mind."} \\
\\ 
\texttt{**8. What is on your right?**  } \\
\texttt{*Answer:* The side of the body, or "the future."} \\
\\ 
\texttt{**9. What is nowhere?**} \\
            }
        }
    }
    & & \\ \\

    \theutterance \stepcounter{utterance}  
    & & & \multicolumn{2}{p{0.3\linewidth}}{
        \cellcolor[rgb]{0.9,0.9,0.9}{
            \makecell[{{p{\linewidth}}}]{
                \texttt{\tiny{[GM$|$GM]}}
                \texttt{Game finished.} \\
            }
        }
    }
    & & \\ \\

\end{supertabular}
}

\end{document}
